% =========================================================
% Analytical Geometry — Notes Index
% =========================================================

% ---------------------------------------------------------
% Breadcrumb
% ---------------------------------------------------------

% ---------------------------------------------------------
% Structural Role
% ---------------------------------------------------------
\begin{tcolorbox}[
  colback=gray!6,
  colframe=gray!40,
  arc=2pt,
  left=8pt, right=8pt, top=6pt, bottom=6pt,
  title={\small\textbf{Structural Role: Analytical Geometry — Algebra Meets Space}},
  fonttitle=\small\bfseries
]
Analytical geometry is the translation layer between algebra and
space.

Every geometric object — a point, line, plane, circle, ellipse —
has an algebraic description, and every algebraic object with
the right structure has a geometric realisation.
This duality is the foundation of the rest of the geometry
sequence.

The key structural tool is the inner product: it encodes both
length and angle in a single algebraic operation.
Orthogonality, projection, and reflection all derive from it,
and the transition to curved geometry replaces the constant
inner product with one that varies from point to point
(the metric tensor in Riemannian geometry).
\end{tcolorbox}
\vspace{1em}

% ---------------------------------------------------------
% Structural Roadmap
% ---------------------------------------------------------
\subsection*{Structural Roadmap}

\begin{center}
\textbf{Definitions $\longrightarrow$ Main Theorems
$\longrightarrow$ Consequences and Structural Insight}
\end{center}

The global progression is:

\begin{enumerate}
  \item Coordinates and the Euclidean inner product
  \item Lengths, angles, and the Cauchy--Schwarz inequality
  \item Lines and planes: parametric, vector, and scalar forms
  \item Projections and orthogonal decompositions
  \item Cross product in $\mathbb{R}^3$: definition and geometry
  \item Conic sections: circles, ellipses, parabolas, hyperbolas
  \item Quadric surfaces: classification via linear algebra
  \item Affine maps and rigid motions
\end{enumerate}

\begin{remark*}[Primary sources]
Axler, \textit{Linear Algebra Done Right}, Chapters~6--7;
Shafarevich \& Remizov, \textit{Linear Algebra and Geometry};
Brannan, Esplen \& Gray, \textit{Geometry}.
\end{remark*}

% ---------------------------------------------------------
% Status
% ---------------------------------------------------------
\vspace{1em}
\begin{tcolorbox}[
  colback=gray!6, colframe=gray!40, arc=2pt,
  left=8pt, right=8pt, top=6pt, bottom=6pt,
  title={\small\textbf{Status: Planned}},
  fonttitle=\small\bfseries
]
\begin{center}\Large\bfseries Coming Soon\end{center}
\vspace{6pt}
\noindent Notes, proofs, and exercises will appear here in a future revision.
\end{tcolorbox}

% Future sections (uncomment as content is written):
% \input{volume-v/geometry/analytical-geometry/notes/notes-inner-product}
% \input{volume-v/geometry/analytical-geometry/notes/notes-lines-planes}
% \input{volume-v/geometry/analytical-geometry/notes/notes-projections}
% \input{volume-v/geometry/analytical-geometry/notes/notes-conics}
% \input{volume-v/geometry/analytical-geometry/notes/notes-affine-maps}
